% !TeX program = lualatex
% ================================================================
% Urdu Tier 3 Vocabulary Version of FactorisingAlgebraicFractionsQuickQuestions
%
% Generated by the server using the following settings:
%
%   language            Urdu (urd)
%   text direction      right to left
%   font                Noto Nastaliq Urdu
%   fallback font       Noto Serif
%   built from          latex/quickQuestions/factorisingAlgebraicFractionsQuickQuestions.tex
%   vocabulary key      latex/quickQuestions/factorisingAlgebraicFractionsQuickQuestions/L2/urd/factorisingAlgebraicFractionsQuickQuestions_urduKey.tex
%   tier 3 vocabulary   green   RGB 0 128 0
%   English text        black   RGB 0 0 0
%   Urdu text           purple  RGB 112 48 160
%   layout macros       version 3
%
% Please don't edit any information in this comment block - the server
% compares it against the current settings and rebuilds this file when
% they differ, so changes here would be overwritten. However, feel free
% to make updates to the LaTeX below if you want to edit it.
% ================================================================

\documentclass[aspectratio=169]{beamer}
% ================================================================
% Copied from the English sheet this was translated from, so that
% anything it sets up for itself still works here
% ================================================================
\usepackage{amsmath}
\usepackage{amssymb}
\usepackage{cancel}
\usepackage{tikz}
\usetikzlibrary{positioning}

% ================================================================
% Quick Questions deck: factorising and algebraic fractions.
%
% Each topic opens with five true/false statements and then runs ten
% quick questions that get gradually harder. Every question gets two
% slides: the question on its own for the mini whiteboards, then the
% same question again with the solution underneath in red.
%
% The two green topics are factorising itself, first into a single
% bracket and then into a double bracket. The two amber topics use
% that factorising to simplify an algebraic fraction, again single
% brackets before double ones. The two red topics do arithmetic with
% algebraic fractions: adding and subtracting, then multiplying and
% dividing.
%
% Where the working can be drawn it is: area rectangles for taking a
% common factor out, the 2 by 2 grid for the double brackets, a
% struck-through common factor for the cancelling, and a flow chart
% for dividing, where turning the second fraction upside down is a
% step of its own. Adding and subtracting is left as written working:
% an algebraic numerator cannot be shaded honestly, so a picture there
% would show a different quantity from the answer.
% ================================================================

\setbeamertemplate{navigation symbols}{}

% A link in the bottom left of every slide back to the contents, so a
% deck this long can be jumped around in during a lesson.
\setbeamertemplate{footline}{%
  \hbox to\paperwidth{%
    \hspace{1em}%
    \hyperlink{contents}{%
      \color{gray}\scriptsize$\leftarrow$~Back to contents}%
    \hfil}%
  \vskip4pt%
}

% Topic divider.
\newcommand{\qtopic}[1]{%
  \section{#1}%
  \begin{frame}
    \vfill
    \begin{center}
      {\usebeamercolor[fg]{structure}\Huge\bfseries #1}
    \end{center}
    \vfill
  \end{frame}%
}

% \qq[<question size>][<solution size>]{<question>}{<solution>}
% The two optional sizes drop the type for the wordier questions and
% the longer solutions, so nothing runs off the slide.
\NewDocumentCommand{\qq}{O{\Huge}O{\LARGE}+m+m}{%
  \begin{frame}
    \vfill
    \begin{center}
      \begin{minipage}{0.92\textwidth}
        \centering #1 #3
      \end{minipage}
    \end{center}
    \vfill
  \end{frame}%
  \begin{frame}
    \vfill
    \begin{center}
      \begin{minipage}{0.92\textwidth}
        \centering #1 #3
        \par\vspace{0.9em}
        {\color{red}#2 #4}
      \end{minipage}
    \end{center}
    \vfill
  \end{frame}%
}

% \tf[<statement size>][<answer size>]{<statement>}{<answer>}
% As \qq, but with a standing "True or false?" heading above the
% statement. The statements hunt for the usual misconceptions, so the
% answer always says why, not just which.
\NewDocumentCommand{\tf}{O{\LARGE}O{\Large}+m+m}{%
  \begin{frame}
    \vfill
    \begin{center}
      \begin{minipage}{0.92\textwidth}
        \centering
        {\usebeamercolor[fg]{structure}\Large\bfseries True or false?}
        \par\vspace{1em}
        #1 #3
      \end{minipage}
    \end{center}
    \vfill
  \end{frame}%
  \begin{frame}
    \vfill
    \begin{center}
      \begin{minipage}{0.92\textwidth}
        \centering
        {\usebeamercolor[fg]{structure}\Large\bfseries True or false?}
        \par\vspace{1em}
        #1 #3
        \par\vspace{0.9em}
        {\color{red}#2 #4}
      \end{minipage}
    \end{center}
    \vfill
  \end{frame}%
}

% Chained working, one equation per row, so that a long simplification
% does not have to be squeezed onto a single line. \dfrac is taller
% than the row strut, so rows carry an explicit skip such as \\[0.95em].
\newcommand{\work}[1]{%
  \begin{tabular}{@{}r@{\;=\;}l@{}}#1\end{tabular}%
}

% \strike{<expression>} --- a factor crossed out because it has been
% divided from the top and the bottom of a fraction.
\newcommand{\strike}[1]{\cancel{#1}}

% ================================================================
% Area pictures for factorising
% ================================================================

% \pairarea{<height>}{<width one>}{<width two>}{<area one>}{<area two>}
% One rectangle of height <height> cut into two parts. The two parts
% have the areas of the two terms being factorised, and the height is
% the factor taken out, so the whole rectangle is the factorised form.
\newcommand{\pairarea}[5]{%
  \begin{tikzpicture}[x=1cm,y=1cm,font=\small]
    \draw[red!60!black,thick,fill=red!8] (0,0) rectangle (3.2,1.5);
    \draw[red!60!black,thick,fill=red!8] (3.2,0) rectangle (5.0,1.5);
    \node[red] at (1.6,0.75) {$#4$};
    \node[red] at (4.1,0.75) {$#5$};
    \node[red,above=2pt] at (1.6,1.5) {$#2$};
    \node[red,above=2pt] at (4.1,1.5) {$#3$};
    \node[red,left=3pt] at (0,0.75) {$#1$};
  \end{tikzpicture}%
}

% \triplearea{<height>}{<w1>}{<w2>}{<w3>}{<a1>}{<a2>}{<a3>}
% The same picture cut into three parts, for a three term expression.
\newcommand{\triplearea}[7]{%
  \begin{tikzpicture}[x=1cm,y=1cm,font=\small]
    \draw[red!60!black,thick,fill=red!8] (0,0) rectangle (2.6,1.5);
    \draw[red!60!black,thick,fill=red!8] (2.6,0) rectangle (4.8,1.5);
    \draw[red!60!black,thick,fill=red!8] (4.8,0) rectangle (6.4,1.5);
    \node[red] at (1.3,0.75) {$#5$};
    \node[red] at (3.7,0.75) {$#6$};
    \node[red] at (5.6,0.75) {$#7$};
    \node[red,above=2pt] at (1.3,1.5) {$#2$};
    \node[red,above=2pt] at (3.7,1.5) {$#3$};
    \node[red,above=2pt] at (5.6,1.5) {$#4$};
    \node[red,left=3pt] at (0,0.75) {$#1$};
  \end{tikzpicture}%
}

% \quadgrid{<w1>}{<w2>}{<h1>}{<h2>}{<a>}{<b>}{<c>}{<d>}
% The 2 by 2 grid for a double bracket. The two widths are the terms
% of one bracket and the two heights the terms of the other, so the
% four cells add to the quadratic and the labels round the outside
% read off the factorised form.
\newcommand{\quadgrid}[8]{%
  \begin{tikzpicture}[x=1cm,y=1cm,font=\small]
    \draw[red!60!black,thick,fill=red!8] (0,0) rectangle (2.8,1.4);
    \draw[red!60!black,thick,fill=red!8] (2.8,0) rectangle (4.6,1.4);
    \draw[red!60!black,thick,fill=red!8] (0,-1.1) rectangle (2.8,0);
    \draw[red!60!black,thick,fill=red!8] (2.8,-1.1) rectangle (4.6,0);
    \node[red] at (1.4,0.7) {$#5$};
    \node[red] at (3.7,0.7) {$#6$};
    \node[red] at (1.4,-0.55) {$#7$};
    \node[red] at (3.7,-0.55) {$#8$};
    \node[red,above=2pt] at (1.4,1.4) {$#1$};
    \node[red,above=2pt] at (3.7,1.4) {$#2$};
    \node[red,left=3pt] at (0,0.7) {$#3$};
    \node[red,left=3pt] at (0,-0.55) {$#4$};
  \end{tikzpicture}%
}

% \flow{<start>}{<first step>}{<middle>}{<second step>}{<end>}
% A flow chart of three boxes joined by two labelled arrows. Used for
% dividing, where turning the second fraction upside down and then
% multiplying are separate steps, done one after the other. The step
% labels wrap to a fixed width, so a label of a few words sits above
% its own arrow instead of running into the boxes either side.
\newcommand{\flow}[5]{%
  \begin{tikzpicture}[
      fbox/.style={draw=red!60!black, rounded corners=2pt, fill=red!8,
                   inner sep=5pt, text=red, font=\Large},
      farr/.style={->, red, thick, shorten >=2pt, shorten <=2pt},
      flab/.style={above, midway, text=red, font=\small,
                   align=center, text width=3.2cm, inner sep=2pt,
                   execute at begin node={\hyphenpenalty=10000
                                          \exhyphenpenalty=10000}}]
    \node[fbox] (a) {$#1$};
    \node[fbox, right=3.4cm of a] (b) {$#3$};
    \node[fbox, right=3.4cm of b] (c) {$#5$};
    \draw[farr] (a) -- node[flab]{#2} (b);
    \draw[farr] (b) -- node[flab]{#4} (c);
  \end{tikzpicture}%
}

% Contents-slide bullets, colour-coded by difficulty band: green for
% the two factorising topics, orange for the two simplifying topics
% and red for the two arithmetic ones. The green is darkened because a
% pure green washes out on a projector.
\setbeamertemplate{section in toc}{%
  \ifnum\inserttocsectionnumber<3
    \textcolor{green!55!black}{$\bullet$}%
  \else
    \ifnum\inserttocsectionnumber<5
      \textcolor{orange}{$\bullet$}%
    \else
      \textcolor{red}{$\bullet$}%
    \fi
  \fi
  ~\inserttocsection%
}

\title{Factorising and Algebraic Fractions: Quick Questions}
\subtitle{Mini whiteboard questions, with solutions}
\author{}
\date{}

% Characters this language's font has not got are borrowed from another,
% rather than being silently dropped - see docs/translated-sheet-latex.md
\directlua{%
  if luaotfload and luaotfload.add_fallback then
    luaotfload.add_fallback("eallatinfallback",
      { "Noto Serif:mode=harf;" })
  else
    tex.error("this luaotfload is too old to register a fallback font")
  end
}
\usepackage[english,bidi=basic]{babel}
\babelprovide[import]{urdu}
\babelfont[urdu]{rm}[Renderer=Harfbuzz, BoldFont={Noto Nastaliq Urdu}, ItalicFont={Noto Nastaliq Urdu}, BoldItalicFont={Noto Nastaliq Urdu}, RawFeature={fallback=eallatinfallback}]{Noto Nastaliq Urdu}
\babelfont[urdu]{sf}[Renderer=Harfbuzz, BoldFont={Noto Nastaliq Urdu}, ItalicFont={Noto Nastaliq Urdu}, BoldItalicFont={Noto Nastaliq Urdu}, RawFeature={fallback=eallatinfallback}]{Noto Nastaliq Urdu}
\newcommand{\ealtext}[1]{\foreignlanguage{urdu}{#1}}
\newcommand{\ealtextblock}[1]{{\raggedleft\foreignlanguage{urdu}{#1}\par}}
% ================================================================
% EAL layout helpers
% ================================================================
\usepackage{xcolor}

\definecolor{ealkeycolour}{RGB}{0,128,0}
\definecolor{ealenglish}{RGB}{0,0,0}
\definecolor{ealtwocolour}{RGB}{112,48,160}

\newsavebox{\ealboxen}
\newsavebox{\ealboxtr}
\newlength{\ealglosswd}
\newlength{\ealglossraise}
\setlength{\ealglossraise}{1.15em}

% a tier 3 word, in English and in translation alike
\newcommand{\ealkey}[1]{\textcolor{ealkeycolour}{#1}}
\newcommand{\ealkeytr}[1]{\textcolor{ealkeycolour}{\ealtext{#1}}}

% One translated word above one English word. Built from kernel
% primitives, never a tabular, which would break wherever the sheet
% loads array, booktabs or tabularx. The box is as wide as the wider of
% the two, so a long translation pushes its neighbours aside instead of
% overlapping them, and it claims its full height so a table row or a
% TikZ node makes room for it.
\newcommand{\ealgloss}[2]{%
  \sbox{\ealboxen}{\ealkey{#1}}%
  \sbox{\ealboxtr}{\small\textcolor{ealtwocolour}{\ealtext{#2}}}%
  \setlength{\ealglosswd}{\wd\ealboxen}%
  \ifdim\wd\ealboxtr>\ealglosswd\setlength{\ealglosswd}{\wd\ealboxtr}\fi
  \leavevmode
  \makebox[\ealglosswd][c]{%
    \makebox[0pt][c]{\usebox{\ealboxen}}%
    \makebox[0pt][c]{\raisebox{\ealglossraise}{\usebox{\ealboxtr}}}%
  }%
}

% opens up line spacing around a block of glosses, so the raised
% translations sit clear of the line above
\newenvironment{ealglossed}{\par\linespread{2.1}\selectfont\ignorespaces}{\par}

% a whole translated sentence above its English counterpart. block level,
% so it does not belong inside a tabular cell
\newcommand{\ealpara}[2]{%
  \par\addvspace{0.5\baselineskip}%
  {\color{ealtwocolour}\ealtextblock{#1}}%
  \penalty200
  {\color{ealenglish}#2\par}%
  \addvspace{0.5\baselineskip}%
}

\begin{document}

\frame{\titlepage}

\begin{frame}[label=contents]{Contents}
  \small
  \tableofcontents
\end{frame}

%================================================================
\qtopic{\ealgloss{Factorising}{تجزیه کرنا} into single \ealgloss{brackets}{بریکٹ}}
%================================================================

\begin{ealglossed}
\tf[\LARGE][\large]{$6x+9=3(2x+3)$}
  {\textbf{True.} The \ealgloss{highest common factor}{سب سے بڑا مشترک عامل} of $6x$ and $9$ is $3$, and
   $3\times 2x=6x$, $3\times 3=9$.
   \par\vspace{0.5em}
   \pairarea{3}{2x}{3}{6x}{9}}
\end{ealglossed}

\begin{ealglossed}
\tf[\Large][\large]{$8x+12$ is fully \ealgloss{factorised}{تجزیه کرنا} as $2(4x+6)$}
  {\textbf{False.} $2$ is a \ealgloss{common factor}{مشترک عامل} of $8x$ and $12$, but it is not
   the \ealgloss{highest common factor}{سب سے بڑا مشترک عامل}, so the \ealgloss{expression}{اظہار} is not fully
   \ealgloss{factorised}{تجزیه کرنا}.\\[0.3em]
   The \ealgloss{highest common factor}{سب سے بڑا مشترک عامل} of $8x$ and $12$ is $4$, so
   $8x+12=4(2x+3)$.}
\end{ealglossed}

\begin{ealglossed}
\tf{$5x+10y=5(x+10y)$}
  {\textbf{False.} Every \ealgloss{term}{رکن} inside the \ealgloss{bracket}{بریکٹ} must be \ealgloss{divided}{تقسیم کرنا} by the
   \ealgloss{common factor}{مشترک عامل}, not only the first \ealgloss{term}{رکن}.\\[0.3em]
   $10y\div 5=2y$, so $5x+10y=5(x+2y)$.}
\end{ealglossed}

\begin{ealglossed}
\tf{$x^{2}+x=x(x)$}
  {\textbf{False.} $x(x)$ is $x^{2}$, so the second \ealgloss{term}{رکن} has been lost.\\[0.3em]
   $x\div x=1$, so $x^{2}+x=x(x+1)$.}
\end{ealglossed}

\begin{ealglossed}
\tf{$3x+3=3(x+1)$}
  {\textbf{True.} The second \ealgloss{term}{رکن} is entirely the \ealgloss{common factor}{مشترک عامل}, and
   $3\div 3=1$, so a $1$ is left inside the \ealgloss{bracket}{بریکٹ}.}
\end{ealglossed}

\begin{ealglossed}
\qq[\Huge][\large]{\ealgloss{Factorise}{تجزیه کرنا} $2x+6$}
  {The \ealgloss{highest common factor}{سب سے بڑا مشترک عامل} of $2x$ and $6$ is $2$.\\[0.3em]
   $2x+6=2(x+3)$
   \par\vspace{0.5em}
   \pairarea{2}{x}{3}{2x}{6}}
\end{ealglossed}

\begin{ealglossed}
\qq[\Huge][\LARGE]{\ealgloss{Factorise}{تجزیه کرنا} $9x+24$}
  {The \ealgloss{highest common factor}{سب سے بڑا مشترک عامل} of $9x$ and $24$ is $3$.\\[0.3em]
   $9x\div 3=3x$ and $24\div 3=8$.\\[0.3em]
   $9x+24=3(3x+8)$}
\end{ealglossed}

\begin{ealglossed}
\qq[\Huge][\LARGE]{\ealgloss{Factorise}{تجزیه کرنا} $6y-15$}
  {The \ealgloss{highest common factor}{سب سے بڑا مشترک عامل} of $6y$ and $15$ is $3$.\\[0.3em]
   $6y-15=3(2y-5)$}
\end{ealglossed}

\begin{ealglossed}
\qq[\Huge][\LARGE]{\ealgloss{Factorise}{تجزیه کرنا} $10a+15b$}
  {The \ealgloss{highest common factor}{سب سے بڑا مشترک عامل} of $10a$ and $15b$ is $5$.\\[0.3em]
   $10a+15b=5(2a+3b)$}
\end{ealglossed}

\begin{ealglossed}
\qq[\Huge][\large]{\ealgloss{Factorise}{تجزیه کرنا} $x^{2}+7x$}
  {The \ealgloss{highest common factor}{سب سے بڑا مشترک عامل} of $x^{2}$ and $7x$ is $x$.\\[0.3em]
   $x^{2}+7x=x(x+7)$
   \par\vspace{0.5em}
   \pairarea{x}{x}{7}{x^{2}}{7x}}
\end{ealglossed}

\begin{ealglossed}
\qq[\Huge][\large]{\ealgloss{Factorise}{تجزیه کرنا} $12x^{2}+8x$}
  {The \ealgloss{highest common factor}{سب سے بڑا مشترک عامل} of $12x^{2}$ and $8x$ is $4x$.\\[0.3em]
   $12x^{2}+8x=4x(3x+2)$
   \par\vspace{0.5em}
   \pairarea{4x}{3x}{2}{12x^{2}}{8x}}
\end{ealglossed}

\begin{ealglossed}
\qq[\Huge][\LARGE]{\ealgloss{Factorise}{تجزیه کرنا} $6x^{3}-9x^{2}$}
  {The \ealgloss{highest common factor}{سب سے بڑا مشترک عامل} of the numbers is $3$ and of the \ealgloss{powers}{قوت} of
   $x$ is $x^{2}$, so the \ealgloss{highest common factor}{سب سے بڑا مشترک عامل} is $3x^{2}$.\\[0.3em]
   $6x^{3}-9x^{2}=3x^{2}(2x-3)$}
\end{ealglossed}

\begin{ealglossed}
\qq[\Huge][\LARGE]{\ealgloss{Factorise}{تجزیه کرنا} $4x^{2}y+6xy^{2}$}
  {Both \ealgloss{terms}{رکن} contain $2$, $x$ and $y$, so the \ealgloss{highest common factor}{سب سے بڑا مشترک عامل} is
   $2xy$.\\[0.3em]
   $4x^{2}y+6xy^{2}=2xy(2x+3y)$}
\end{ealglossed}

\begin{ealglossed}
\qq[\Huge][\large]{\ealgloss{Factorise}{تجزیه کرنا} $8x^{3}+12x^{2}+20x$}
  {The \ealgloss{highest common factor}{سب سے بڑا مشترک عامل} of $8x^{3}$, $12x^{2}$ and $20x$ is $4x$.\\[0.3em]
   $8x^{3}+12x^{2}+20x=4x(2x^{2}+3x+5)$
   \par\vspace{0.5em}
   \triplearea{4x}{2x^{2}}{3x}{5}{8x^{3}}{12x^{2}}{20x}}
\end{ealglossed}

\begin{ealglossed}
\qq[\LARGE][\large]{\ealgloss{Factorise}{تجزیه کرنا} $5(x+1)+x(x+1)$}
  {The \ealgloss{bracket}{بریکٹ} $(x+1)$ is a \ealgloss{common factor}{مشترک عامل} of both \ealgloss{terms}{رکن}.\\[0.3em]
   $5(x+1)+x(x+1)=(x+1)(x+5)$
   \par\vspace{0.5em}
   \pairarea{x+1}{x}{5}{x(x+1)}{5(x+1)}}
\end{ealglossed}

%================================================================
\qtopic{\ealgloss{Factorising}{تجزیه کرنا} into double \ealgloss{brackets}{بریکٹ}}
%================================================================

\begin{ealglossed}
\tf[\LARGE][\large]{$x^{2}-9=(x-3)(x-3)$}
  {\textbf{False.} $(x-3)(x-3)$ \ealgloss{expands}{توسیع کرنا} to $x^{2}-6x+9$.\\[0.3em]
   The \ealgloss{signs}{علامت} must be opposite: $x^{2}-9=(x+3)(x-3)$.
   \par\vspace{0.4em}
   \quadgrid{x}{3}{x}{-3}{x^2}{3x}{-3x}{-9}}
\end{ealglossed}

\begin{ealglossed}
\tf{$x^{2}+9=(x+3)(x+3)$}
  {\textbf{False.} $(x+3)(x+3)$ \ealgloss{expands}{توسیع کرنا} to $x^{2}+6x+9$.\\[0.3em]
   $x^{2}+9$ is a \ealgloss{sum}{حاصل جمع} of two \ealgloss{squares}{مربع}, and it does not \ealgloss{factorise}{تجزیه کرنا}.}
\end{ealglossed}

\tf{$x^{2}+2x-15=(x+5)(x-3)$}
  {\textbf{True.} $5\times(-3)=-15$ and $5+(-3)=2$.}

\begin{ealglossed}
\tf{$x^{2}-2x-8=(x+4)(x-2)$}
  {\textbf{False.} $(x+4)(x-2)$ \ealgloss{expands}{توسیع کرنا} to $x^{2}+2x-8$, so the \ealgloss{signs}{علامت} are
   the wrong way round.\\[0.3em]
   $x^{2}-2x-8=(x-4)(x+2)$}
\end{ealglossed}

\begin{ealglossed}
\tf[\Large][\large]{To \ealgloss{factorise}{تجزیه کرنا} $x^{2}+7x+12$ we need two numbers\\[0.3em]
   that \ealgloss{multiply}{ضرب دینا} to $12$ and \ealgloss{add}{جمع کرنا} to $7$}
  {\textbf{True.} The two numbers \ealgloss{multiply}{ضرب دینا} to the \ealgloss{constant term}{مستقل حد} and \ealgloss{add}{جمع کرنا}
   to the \ealgloss{coefficient}{عددی سر} of $x$.\\[0.3em]
   The numbers are $3$ and $4$, so $x^{2}+7x+12=(x+3)(x+4)$.}
\end{ealglossed}

\begin{ealglossed}
\qq[\Huge][\large]{\ealgloss{Factorise}{تجزیه کرنا} $x^{2}+5x+6$}
  {Two numbers that \ealgloss{multiply}{ضرب دینا} to $6$ and \ealgloss{add}{جمع کرنا} to $5$: $2$ and $3$.\\[0.3em]
   $x^{2}+5x+6=(x+2)(x+3)$
   \par\vspace{0.5em}
   \quadgrid{x}{2}{x}{3}{x^2}{2x}{3x}{6}}
\end{ealglossed}

\begin{ealglossed}
\qq[\Huge][\LARGE]{\ealgloss{Factorise}{تجزیه کرنا} $x^{2}+9x+20$}
  {Two numbers that \ealgloss{multiply}{ضرب دینا} to $20$ and \ealgloss{add}{جمع کرنا} to $9$: $4$ and $5$.\\[0.3em]
   $x^{2}+9x+20=(x+4)(x+5)$}
\end{ealglossed}

\begin{ealglossed}
\qq[\Huge][\large]{\ealgloss{Factorise}{تجزیه کرنا} $x^{2}+3x-10$}
  {Two numbers that \ealgloss{multiply}{ضرب دینا} to $-10$ and \ealgloss{add}{جمع کرنا} to $3$: $5$ and $-2$.\\[0.3em]
   $x^{2}+3x-10=(x+5)(x-2)$
   \par\vspace{0.5em}
   \quadgrid{x}{5}{x}{-2}{x^2}{5x}{-2x}{-10}}
\end{ealglossed}

\begin{ealglossed}
\qq[\Huge][\LARGE]{\ealgloss{Factorise}{تجزیه کرنا} $x^{2}-7x+12$}
  {The \ealgloss{constant term}{مستقل حد} is positive and the \ealgloss{coefficient}{عددی سر} of $x$ is negative,
   so both numbers are negative: $-3$ and $-4$.\\[0.3em]
   $x^{2}-7x+12=(x-3)(x-4)$}
\end{ealglossed}

\begin{ealglossed}
\qq[\Huge][\LARGE]{\ealgloss{Factorise}{تجزیه کرنا} $x^{2}-2x-15$}
  {Two numbers that \ealgloss{multiply}{ضرب دینا} to $-15$ and \ealgloss{add}{جمع کرنا} to $-2$: $-5$ and $3$.\\[0.3em]
   $x^{2}-2x-15=(x-5)(x+3)$}
\end{ealglossed}

\begin{ealglossed}
\qq[\Huge][\LARGE]{\ealgloss{Factorise}{تجزیه کرنا} $x^{2}-49$}
  {This is a \ealgloss{difference of two squares}{دو مربعوں کا فرق}, with $x^{2}=(x)^{2}$ and
   $49=7^{2}$.\\[0.3em]
   $x^{2}-49=(x+7)(x-7)$}
\end{ealglossed}

\begin{ealglossed}
\qq[\Huge][\large]{\ealgloss{Factorise}{تجزیه کرنا} $2x^{2}+7x+3$}
  {The $2x^{2}$ comes from $2x\times x$, and the $3$ from $1\times 3$.\\[0.3em]
   $2x^{2}+7x+3=(2x+1)(x+3)$
   \par\vspace{0.5em}
   \quadgrid{2x}{1}{x}{3}{2x^2}{x}{6x}{3}}
\end{ealglossed}

\begin{ealglossed}
\qq[\Huge][\large]{\ealgloss{Factorise}{تجزیه کرنا} $2x^{2}+10x+12$}
  {Take the \ealgloss{common factor}{مشترک عامل} of $2$ out first.\\[0.3em]
   $2x^{2}+10x+12=2(x^{2}+5x+6)$\\[0.3em]
   $x^{2}+5x+6=(x+2)(x+3)$, so $2x^{2}+10x+12=2(x+2)(x+3)$.
   \par\vspace{0.5em}
   \triplearea{2}{x^{2}}{5x}{6}{2x^{2}}{10x}{12}}
\end{ealglossed}

\begin{ealglossed}
\qq[\Huge][\LARGE]{\ealgloss{Factorise}{تجزیه کرنا} $9x^{2}-16$}
  {This is a \ealgloss{difference of two squares}{دو مربعوں کا فرق}, with $9x^{2}=(3x)^{2}$ and
   $16=4^{2}$.\\[0.3em]
   $9x^{2}-16=(3x+4)(3x-4)$}
\end{ealglossed}

\begin{ealglossed}
\qq[\Huge][\large]{\ealgloss{Factorise}{تجزیه کرنا} $6x^{2}-x-12$}
  {The $6x^{2}$ comes from $3x\times 2x$, and the $-12$ from
   $4\times(-3)$.\\[0.3em]
   $6x^{2}-x-12=(3x+4)(2x-3)$
   \par\vspace{0.5em}
   \quadgrid{3x}{4}{2x}{-3}{6x^2}{8x}{-9x}{-12}}
\end{ealglossed}

%================================================================
\qtopic{\ealgloss{Simplifying}{سادہ کرنا} \ealgloss{algebraic}{الجبری} \ealgloss{fractions}{کسر} with single \ealgloss{brackets}{بریکٹ}}
%================================================================

\begin{ealglossed}
\tf[\LARGE][\large]{$\dfrac{x+2}{2}=x$}
  {\textbf{False.} The $2$ on the bottom \ealgloss{divides}{تقسیم کرنا} the whole \ealgloss{numerator}{صورت},
   not the \ealgloss{term}{رکن} $2$ on its own.\\[0.3em]
   $x+2$ has no \ealgloss{common factor}{مشترک عامل}, so $\dfrac{x+2}{2}$ does not
   \ealgloss{simplify}{سادہ کرنا}.\\[0.3em]
   When $x=4$: $\dfrac{4+2}{2}=3$, not $4$.}
\end{ealglossed}

\begin{ealglossed}
\tf[\LARGE][\large]{$\dfrac{3x+6}{3}=x+2$}
  {\textbf{True.} $3$ is a \ealgloss{factor}{عامل} of the whole \ealgloss{numerator}{صورت}:
   $3x+6=3(x+2)$.\\[0.3em]
   $\dfrac{\strike{3}(x+2)}{\strike{3}}=x+2$}
\end{ealglossed}

\begin{ealglossed}
\tf[\LARGE][\large]{$\dfrac{x+3}{x}=3$}
  {\textbf{False.} $x$ is a \ealgloss{term}{رکن} of the \ealgloss{numerator}{صورت}, not a \ealgloss{factor}{عامل} of the
   whole \ealgloss{numerator}{صورت}.\\[0.3em]
   $\dfrac{x+3}{x}$ does not \ealgloss{simplify}{سادہ کرنا}.\\[0.3em]
   When $x=1$: $\dfrac{1+3}{1}=4$, not $3$.}
\end{ealglossed}

\begin{ealglossed}
\tf[\LARGE][\large]{$\dfrac{6x+9}{3}=2x+9$}
  {\textbf{False.} Every \ealgloss{term}{رکن} of the \ealgloss{numerator}{صورت} is \ealgloss{divided}{تقسیم کرنا} by $3$, not
   the first \ealgloss{term}{رکن} on its own.\\[0.3em]
   $6x+9=3(2x+3)$, so the \ealgloss{fraction}{کسر} \ealgloss{simplifies}{سادہ کرنا} to $2x+3$.}
\end{ealglossed}

\begin{ealglossed}
\tf[\LARGE][\large]{$\dfrac{5x+10}{5x}=\dfrac{x+2}{x}$}
  {\textbf{True.} $5$ is a \ealgloss{common factor}{مشترک عامل} of the whole \ealgloss{numerator}{صورت} and the
   whole \ealgloss{denominator}{مخرج}.\\[0.3em]
   $\dfrac{\strike{5}(x+2)}{\strike{5}x}=\dfrac{x+2}{x}$}
\end{ealglossed}

\begin{ealglossed}
\qq[\LARGE][\normalsize]{\ealgloss{Simplify}{سادہ کرنا} $\dfrac{2x+10}{2}$}
  {\ealgloss{Factorise}{تجزیه کرنا} the \ealgloss{numerator}{صورت}.\\[0.3em]
   $\dfrac{2x+10}{2}=\dfrac{\strike{2}(x+5)}{\strike{2}}=x+5$
   \par\vspace{0.5em}
   \pairarea{2}{x}{5}{2x}{10}}
\end{ealglossed}

\begin{ealglossed}
\qq[\LARGE][\large]{\ealgloss{Simplify}{سادہ کرنا} $\dfrac{5x-15}{5}$}
  {The \ealgloss{highest common factor}{سب سے بڑا مشترک عامل} of the \ealgloss{numerator}{صورت} is $5$.\\[0.3em]
   $\dfrac{5(x-3)}{5}=\dfrac{\strike{5}(x-3)}{\strike{5}}=x-3$}
\end{ealglossed}

\begin{ealglossed}
\qq[\LARGE][\normalsize]{\ealgloss{Simplify}{سادہ کرنا} $\dfrac{2x+8}{x+4}$}
  {\ealgloss{Factorise}{تجزیه کرنا} the \ealgloss{numerator}{صورت}.\\[0.3em]
   $\dfrac{2x+8}{x+4}=\dfrac{2\strike{(x+4)}}{\strike{(x+4)}}=2$
   \par\vspace{0.5em}
   \pairarea{2}{x}{4}{2x}{8}}
\end{ealglossed}

\begin{ealglossed}
\qq[\LARGE][\large]{\ealgloss{Simplify}{سادہ کرنا} $\dfrac{3x-12}{x-4}$}
  {\ealgloss{Factorise}{تجزیه کرنا} the \ealgloss{numerator}{صورت}: $3x-12=3(x-4)$.\\[0.3em]
   $\dfrac{3\strike{(x-4)}}{\strike{(x-4)}}=3$}
\end{ealglossed}

\begin{ealglossed}
\qq[\LARGE][\normalsize]{\ealgloss{Simplify}{سادہ کرنا} $\dfrac{x^{2}+3x}{x}$}
  {The \ealgloss{highest common factor}{سب سے بڑا مشترک عامل} of the \ealgloss{numerator}{صورت} is $x$.\\[0.3em]
   $\dfrac{x(x+3)}{x}=\dfrac{\strike{x}(x+3)}{\strike{x}}=x+3$
   \par\vspace{0.5em}
   \pairarea{x}{x}{3}{x^{2}}{3x}}
\end{ealglossed}

\begin{ealglossed}
\qq[\LARGE][\normalsize]{\ealgloss{Simplify}{سادہ کرنا} $\dfrac{3x^{2}+6x+9}{3}$}
  {$3$ is a \ealgloss{factor}{عامل} of all three \ealgloss{terms}{رکن} of the \ealgloss{numerator}{صورت}.\\[0.3em]
   $\dfrac{\strike{3}(x^{2}+2x+3)}{\strike{3}}=x^{2}+2x+3$
   \par\vspace{0.5em}
   \triplearea{3}{x^{2}}{2x}{3}{3x^{2}}{6x}{9}}
\end{ealglossed}

\begin{ealglossed}
\qq[\LARGE][\normalsize]{\ealgloss{Simplify}{سادہ کرنا} $\dfrac{4x^{2}+6x}{2x}$}
  {The \ealgloss{highest common factor}{سب سے بڑا مشترک عامل} of the \ealgloss{numerator}{صورت} is $2x$.\\[0.3em]
   $\dfrac{2x(2x+3)}{2x}=\dfrac{\strike{2x}(2x+3)}{\strike{2x}}=2x+3$
   \par\vspace{0.5em}
   \pairarea{2x}{2x}{3}{4x^{2}}{6x}}
\end{ealglossed}

\begin{ealglossed}
\qq[\LARGE][\large]{\ealgloss{Simplify}{سادہ کرنا} $\dfrac{4x+12}{6x+18}$}
  {\ealgloss{Factorise}{تجزیه کرنا} the \ealgloss{numerator}{صورت} and the \ealgloss{denominator}{مخرج}.
   \par\vspace{0.4em}
   \work{$\dfrac{4x+12}{6x+18}$ & $\dfrac{4(x+3)}{6(x+3)}$\\[0.95em]
         $\dfrac{4\strike{(x+3)}}{6\strike{(x+3)}}$ & $\dfrac{4}{6}=\dfrac{2}{3}$}}
\end{ealglossed}

\begin{ealglossed}
\qq[\LARGE][\large]{\ealgloss{Simplify}{سادہ کرنا} $\dfrac{10}{5x+20}$}
  {\ealgloss{Factorise}{تجزیه کرنا} the \ealgloss{denominator}{مخرج}: $5x+20=5(x+4)$.\\[0.3em]
   $\dfrac{\strike{5}\times 2}{\strike{5}(x+4)}=\dfrac{2}{x+4}$}
\end{ealglossed}

\begin{ealglossed}
\qq[\LARGE][\large]{\ealgloss{Simplify}{سادہ کرنا} $\dfrac{6x^{2}+9x}{4x^{2}+6x}$}
  {The \ealgloss{highest common factor}{سب سے بڑا مشترک عامل} is $3x$ on the top and $2x$ on the bottom.
   \par\vspace{0.4em}
   \work{$\dfrac{6x^{2}+9x}{4x^{2}+6x}$ & $\dfrac{3x(2x+3)}{2x(2x+3)}$\\[0.95em]
         $\dfrac{3\strike{x}\strike{(2x+3)}}{2\strike{x}\strike{(2x+3)}}$ & $\dfrac{3}{2}$}}
\end{ealglossed}

%================================================================
\qtopic{\ealgloss{Simplifying}{سادہ کرنا} \ealgloss{algebraic}{الجبری} \ealgloss{fractions}{کسر} with double \ealgloss{brackets}{بریکٹ}}
%================================================================

\begin{ealglossed}
\tf[\LARGE][\large]{$\dfrac{x^{2}+5x+6}{x+2}=x+3$}
  {\textbf{True.} $x^{2}+5x+6=(x+2)(x+3)$, so $(x+2)$ is a \ealgloss{factor}{عامل} of
   the whole \ealgloss{numerator}{صورت} and \ealgloss{cancels}{حذف کرنا}.\\[0.3em]
   $\dfrac{\strike{(x+2)}(x+3)}{\strike{(x+2)}}=x+3$}
\end{ealglossed}

\begin{ealglossed}
\tf[\LARGE][\large]{$\dfrac{x^{2}+7x}{x^{2}+2x}=\dfrac{7}{2}$}
  {\textbf{False.} $x^{2}$ is a \ealgloss{term}{رکن}, not a \ealgloss{factor}{عامل} of the whole
   \ealgloss{numerator}{صورت}, so it cannot be \ealgloss{cancelled}{حذف کرنا}.\\[0.3em]
   $\dfrac{x(x+7)}{x(x+2)}=\dfrac{x+7}{x+2}$}
\end{ealglossed}

\begin{ealglossed}
\tf[\LARGE][\large]{$\dfrac{x-3}{3-x}=1$}
  {\textbf{False.} The two \ealgloss{brackets}{بریکٹ} are not the same:
   $3-x=-(x-3)$.\\[0.3em]
   $\dfrac{x-3}{-(x-3)}=-1$\\[0.3em]
   When $x=5$: $\dfrac{2}{-2}=-1$.}
\end{ealglossed}

\begin{ealglossed}
\tf[\LARGE][\large]{$\dfrac{x^{2}-25}{x-5}=x+5$}
  {\textbf{True.} The \ealgloss{numerator}{صورت} is a \ealgloss{difference of two squares}{دو مربعوں کا فرق}:
   $x^{2}-25=(x-5)(x+5)$.\\[0.3em]
   $\dfrac{\strike{(x-5)}(x+5)}{\strike{(x-5)}}=x+5$}
\end{ealglossed}

\begin{ealglossed}
\tf[\LARGE][\large]{$\dfrac{x^{2}-16}{x^{2}+8x+16}=1$}
  {\textbf{False.} $(x-4)$ and $(x+4)$ are different \ealgloss{brackets}{بریکٹ} and do
   not \ealgloss{cancel}{حذف کرنا}.\\[0.3em]
   $\dfrac{(x-4)(x+4)}{(x+4)(x+4)}=\dfrac{x-4}{x+4}$}
\end{ealglossed}

\begin{ealglossed}
\qq[\LARGE][\large]{\ealgloss{Simplify}{سادہ کرنا} $\dfrac{x^{2}+8x+15}{x+5}$}
  {\ealgloss{Factorise}{تجزیه کرنا} the \ealgloss{numerator}{صورت}. Two numbers with a \ealgloss{sum}{حاصل جمع} of $8$ and a \ealgloss{product}{حاصل ضرب}
   of $15$: $3$ and $5$.\\[0.3em]
   $\dfrac{(x+3)\strike{(x+5)}}{\strike{(x+5)}}=x+3$}
\end{ealglossed}

\begin{ealglossed}
\qq[\LARGE][\large]{\ealgloss{Simplify}{سادہ کرنا} $\dfrac{x^{2}+7x+12}{x+3}$}
  {Two numbers with a \ealgloss{sum}{حاصل جمع} of $7$ and a \ealgloss{product}{حاصل ضرب} of $12$: $3$ and $4$.\\[0.3em]
   $\dfrac{\strike{(x+3)}(x+4)}{\strike{(x+3)}}=x+4$}
\end{ealglossed}

\begin{ealglossed}
\qq[\LARGE][\large]{\ealgloss{Simplify}{سادہ کرنا} $\dfrac{x^{2}-36}{x+6}$}
  {The \ealgloss{numerator}{صورت} is a \ealgloss{difference of two squares}{دو مربعوں کا فرق}:
   $x^{2}-36=(x+6)(x-6)$.\\[0.3em]
   $\dfrac{\strike{(x+6)}(x-6)}{\strike{(x+6)}}=x-6$}
\end{ealglossed}

\begin{ealglossed}
\qq[\LARGE][\normalsize]{\ealgloss{Simplify}{سادہ کرنا} $\dfrac{x^{2}+2x-8}{x-2}$}
  {Two numbers with a \ealgloss{sum}{حاصل جمع} of $2$ and a \ealgloss{product}{حاصل ضرب} of $-8$: $4$ and $-2$.\\[0.3em]
   $\dfrac{(x+4)\strike{(x-2)}}{\strike{(x-2)}}=x+4$
   \par\vspace{0.5em}
   \quadgrid{x}{4}{x}{-2}{x^{2}}{4x}{-2x}{-8}}
\end{ealglossed}

\begin{ealglossed}
\qq[\LARGE][\large]{\ealgloss{Simplify}{سادہ کرنا} $\dfrac{x^{2}+6x+8}{x^{2}+5x+4}$}
  {\ealgloss{Factorise}{تجزیه کرنا} the \ealgloss{numerator}{صورت} and the \ealgloss{denominator}{مخرج}.
   \par\vspace{0.4em}
   \work{$\dfrac{x^{2}+6x+8}{x^{2}+5x+4}$ & $\dfrac{(x+2)(x+4)}{(x+1)(x+4)}$\\[0.95em]
         $\dfrac{(x+2)\strike{(x+4)}}{(x+1)\strike{(x+4)}}$ & $\dfrac{x+2}{x+1}$}}
\end{ealglossed}

\begin{ealglossed}
\qq[\LARGE][\large]{\ealgloss{Simplify}{سادہ کرنا} $\dfrac{x^{2}-x-6}{x^{2}-9}$}
  {The \ealgloss{denominator}{مخرج} is a \ealgloss{difference of two squares}{دو مربعوں کا فرق}.
   \par\vspace{0.4em}
   \work{$\dfrac{x^{2}-x-6}{x^{2}-9}$ & $\dfrac{(x-3)(x+2)}{(x-3)(x+3)}$\\[0.95em]
         $\dfrac{\strike{(x-3)}(x+2)}{\strike{(x-3)}(x+3)}$ & $\dfrac{x+2}{x+3}$}}
\end{ealglossed}

\begin{ealglossed}
\qq[\LARGE][\large]{\ealgloss{Simplify}{سادہ کرنا} $\dfrac{3x+6}{x^{2}+5x+6}$}
  {The \ealgloss{numerator}{صورت} has a \ealgloss{common factor}{مشترک عامل} of $3$; the \ealgloss{denominator}{مخرج}
   \ealgloss{factorises}{تجزیه کرنا} into two \ealgloss{brackets}{بریکٹ}.
   \par\vspace{0.4em}
   \work{$\dfrac{3x+6}{x^{2}+5x+6}$ & $\dfrac{3(x+2)}{(x+2)(x+3)}$\\[0.95em]
         $\dfrac{3\strike{(x+2)}}{\strike{(x+2)}(x+3)}$ & $\dfrac{3}{x+3}$}}
\end{ealglossed}

\begin{ealglossed}
\qq[\LARGE][\normalsize]{\ealgloss{Simplify}{سادہ کرنا} $\dfrac{2x^{2}+7x+3}{x+3}$}
  {The \ealgloss{coefficient}{عددی سر} of $x^{2}$ is $2$, so one \ealgloss{bracket}{بریکٹ} starts with $2x$.\\[0.3em]
   $\dfrac{(2x+1)\strike{(x+3)}}{\strike{(x+3)}}=2x+1$
   \par\vspace{0.5em}
   \quadgrid{2x}{1}{x}{3}{2x^{2}}{x}{6x}{3}}
\end{ealglossed}

\begin{ealglossed}
\qq[\LARGE][\large]{\ealgloss{Simplify}{سادہ کرنا} $\dfrac{2x^{2}-8}{x^{2}+4x+4}$}
  {Take the \ealgloss{common factor}{مشترک عامل} of $2$ out of the \ealgloss{numerator}{صورت} first.
   \par\vspace{0.4em}
   \work{$\dfrac{2x^{2}-8}{x^{2}+4x+4}$ & $\dfrac{2(x^{2}-4)}{(x+2)(x+2)}$\\[0.95em]
         $\dfrac{2(x^{2}-4)}{(x+2)(x+2)}$ & $\dfrac{2(x-2)\strike{(x+2)}}{(x+2)\strike{(x+2)}}$\\[0.95em]
         $\dfrac{2x^{2}-8}{x^{2}+4x+4}$ & $\dfrac{2(x-2)}{x+2}$}}
\end{ealglossed}

\begin{ealglossed}
\qq[\LARGE][\large]{\ealgloss{Simplify}{سادہ کرنا} $\dfrac{6x^{2}+x-2}{4x^{2}-1}$}
  {The \ealgloss{denominator}{مخرج} is a \ealgloss{difference of two squares}{دو مربعوں کا فرق}.
   \par\vspace{0.4em}
   \work{$6x^{2}+x-2$ & $(3x+2)(2x-1)$\\[0.95em]
         $4x^{2}-1$ & $(2x-1)(2x+1)$\\[0.95em]
         $\dfrac{(3x+2)\strike{(2x-1)}}{(2x+1)\strike{(2x-1)}}$ & $\dfrac{3x+2}{2x+1}$}}
\end{ealglossed}

%================================================================
\qtopic{\ealgloss{Adding}{جمع کرنا} and \ealgloss{subtracting}{تفریق کرنا} \ealgloss{algebraic}{الجبری} \ealgloss{fractions}{کسر}}
%================================================================

\begin{ealglossed}
\tf[\LARGE][\Large]{$\dfrac{2}{x}+\dfrac{3}{x}=\dfrac{5}{2x}$}
  {\textbf{False.} The \ealgloss{denominators}{مخرج} are already the same, so only the
   \ealgloss{numerators}{صورت} are \ealgloss{added}{جمع کرنا}.\\[0.3em]
   $\dfrac{2}{x}+\dfrac{3}{x}=\dfrac{5}{x}$}
\end{ealglossed}

\begin{ealglossed}
\tf[\LARGE][\Large]{$\dfrac{x}{3}+\dfrac{x}{4}=\dfrac{7x}{12}$}
  {\textbf{True.} The \ealgloss{common denominator}{مشترک مخرج} is $12$.\\[0.3em]
   $\dfrac{4x}{12}+\dfrac{3x}{12}=\dfrac{7x}{12}$}
\end{ealglossed}

\begin{ealglossed}
\tf[\LARGE][\large]{$\dfrac{1}{x}+\dfrac{1}{y}=\dfrac{1}{x+y}$}
  {\textbf{False.} Test it with $x=2$ and $y=3$: the left side is
   $\dfrac{5}{6}$ and the right side is $\dfrac{1}{5}$.\\[0.3em]
   The \ealgloss{common denominator}{مشترک مخرج} is $xy$, so the answer is $\dfrac{x+y}{xy}$.}
\end{ealglossed}

\begin{ealglossed}
\tf[\Large][\large]{$\dfrac{4}{x}-\dfrac{x-3}{x}=\dfrac{4-x-3}{x}$}
  {\textbf{False.} The minus \ealgloss{sign}{علامت} applies to the whole second \ealgloss{numerator}{صورت},
   so the $-3$ becomes $+3$.\\[0.3em]
   $\dfrac{4-(x-3)}{x}=\dfrac{4-x+3}{x}=\dfrac{7-x}{x}$}
\end{ealglossed}

\begin{ealglossed}
\tf[\Large][\large]{The \ealgloss{common denominator}{مشترک مخرج} of $\dfrac{1}{x+1}$ and
   $\dfrac{2}{x+3}$ is $(x+1)(x+3)$}
  {\textbf{True.} The \ealgloss{denominators}{مخرج} are \ealgloss{multiplied}{ضرب دینا}, not \ealgloss{added}{جمع کرنا}.\\[0.3em]
   The \ealgloss{common denominator}{مشترک مخرج} is $(x+1)(x+3)$, not $x+4$.}
\end{ealglossed}

\begin{ealglossed}
\qq[\LARGE][\large]{Write as a single \ealgloss{fraction}{کسر}\\[0.4em]
   $\dfrac{3}{x}+\dfrac{4}{x}$}
  {The \ealgloss{denominators}{مخرج} are the same, so the \ealgloss{numerators}{صورت} are \ealgloss{added}{جمع کرنا}.\\[0.4em]
   $\dfrac{3}{x}+\dfrac{4}{x}=\dfrac{7}{x}$}
\end{ealglossed}

\begin{ealglossed}
\qq[\LARGE][\large]{Write as a single \ealgloss{fraction}{کسر}\\[0.4em]
   $\dfrac{x}{2}+\dfrac{x}{5}$}
  {The \ealgloss{common denominator}{مشترک مخرج} is $10$.
   \par\vspace{0.4em}
   \work{$\dfrac{x}{2}+\dfrac{x}{5}$ & $\dfrac{5x}{10}+\dfrac{2x}{10}$\\[0.95em]
         $\dfrac{5x}{10}+\dfrac{2x}{10}$ & $\dfrac{7x}{10}$}}
\end{ealglossed}

\begin{ealglossed}
\qq[\LARGE][\large]{Write as a single \ealgloss{fraction}{کسر}\\[0.4em]
   $\dfrac{x}{4}+\dfrac{3x}{8}$}
  {$8$ is a \ealgloss{multiple}{مضاعف} of $4$, so the \ealgloss{common denominator}{مشترک مخرج} is $8$.
   \par\vspace{0.4em}
   \work{$\dfrac{x}{4}+\dfrac{3x}{8}$ & $\dfrac{2x}{8}+\dfrac{3x}{8}$\\[0.95em]
         $\dfrac{2x}{8}+\dfrac{3x}{8}$ & $\dfrac{5x}{8}$}}
\end{ealglossed}

\begin{ealglossed}
\qq[\LARGE][\large]{Write as a single \ealgloss{fraction}{کسر}\\[0.4em]
   $\dfrac{2x}{5}-\dfrac{x}{3}$}
  {The \ealgloss{common denominator}{مشترک مخرج} is $15$.
   \par\vspace{0.4em}
   \work{$\dfrac{2x}{5}-\dfrac{x}{3}$ & $\dfrac{6x}{15}-\dfrac{5x}{15}$\\[0.95em]
         $\dfrac{6x}{15}-\dfrac{5x}{15}$ & $\dfrac{x}{15}$}}
\end{ealglossed}

\begin{ealglossed}
\qq[\LARGE][\large]{Write as a single \ealgloss{fraction}{کسر}\\[0.4em]
   $\dfrac{x+1}{2}+\dfrac{x+2}{3}$}
  {The \ealgloss{common denominator}{مشترک مخرج} is $6$. Each \ealgloss{numerator}{صورت} is \ealgloss{multiplied}{ضرب دینا} by the
   same \ealgloss{factor}{عامل} as its \ealgloss{denominator}{مخرج}.
   \par\vspace{0.4em}
   \work{$\dfrac{x+1}{2}+\dfrac{x+2}{3}$ & $\dfrac{3(x+1)+2(x+2)}{6}$\\[0.95em]
         $\dfrac{3x+3+2x+4}{6}$ & $\dfrac{5x+7}{6}$}}
\end{ealglossed}

\begin{ealglossed}
\qq[\LARGE][\large]{Write as a single \ealgloss{fraction}{کسر}\\[0.4em]
   $\dfrac{2x+3}{4}-\dfrac{x-1}{6}$}
  {The \ealgloss{common denominator}{مشترک مخرج} is $12$. The minus \ealgloss{sign}{علامت} applies to the whole
   second \ealgloss{numerator}{صورت}.
   \par\vspace{0.4em}
   \work{$\dfrac{2x+3}{4}-\dfrac{x-1}{6}$ & $\dfrac{3(2x+3)-2(x-1)}{12}$\\[0.95em]
         $\dfrac{6x+9-2x+2}{12}$ & $\dfrac{4x+11}{12}$}}
\end{ealglossed}

\begin{ealglossed}
\qq[\LARGE][\large]{Write as a single \ealgloss{fraction}{کسر}\\[0.4em]
   $\dfrac{2}{x}+\dfrac{3}{y}$}
  {The \ealgloss{common denominator}{مشترک مخرج} is $xy$.
   \par\vspace{0.4em}
   \work{$\dfrac{2}{x}+\dfrac{3}{y}$ & $\dfrac{2y}{xy}+\dfrac{3x}{xy}$\\[0.95em]
         $\dfrac{2y}{xy}+\dfrac{3x}{xy}$ & $\dfrac{3x+2y}{xy}$}}
\end{ealglossed}

\begin{ealglossed}
\qq[\LARGE][\large]{Write as a single \ealgloss{fraction}{کسر}\\[0.4em]
   $\dfrac{1}{x+1}+\dfrac{2}{x+3}$}
  {The \ealgloss{common denominator}{مشترک مخرج} is $(x+1)(x+3)$.
   \par\vspace{0.4em}
   \work{$\dfrac{1}{x+1}+\dfrac{2}{x+3}$ & $\dfrac{(x+3)+2(x+1)}{(x+1)(x+3)}$\\[0.95em]
         $\dfrac{x+3+2x+2}{(x+1)(x+3)}$ & $\dfrac{3x+5}{(x+1)(x+3)}$}}
\end{ealglossed}

\begin{ealglossed}
\qq[\LARGE][\large]{Write as a single \ealgloss{fraction}{کسر}\\[0.4em]
   $\dfrac{3}{x-1}-\dfrac{2}{x+2}$}
  {The \ealgloss{common denominator}{مشترک مخرج} is $(x-1)(x+2)$. The minus \ealgloss{sign}{علامت} applies to the
   whole second \ealgloss{numerator}{صورت}.
   \par\vspace{0.4em}
   \work{$\dfrac{3}{x-1}-\dfrac{2}{x+2}$ & $\dfrac{3(x+2)-2(x-1)}{(x-1)(x+2)}$\\[0.95em]
         $\dfrac{3x+6-2x+2}{(x-1)(x+2)}$ & $\dfrac{x+8}{(x-1)(x+2)}$}}
\end{ealglossed}

\begin{ealglossed}
\qq[\LARGE][\small]{Write as a single \ealgloss{fraction}{کسر}\\[0.4em]
   $\dfrac{1}{x-2}+\dfrac{3}{x^{2}-4}$}
  {\ealgloss{Factorise}{تجزیه کرنا} the second \ealgloss{denominator}{مخرج}, then write both \ealgloss{fractions}{کسر} over
   $(x-2)(x+2)$.
   \par\vspace{0.4em}
   \work{$x^{2}-4$ & $(x-2)(x+2)$\\[1.3em]
         $\dfrac{1}{x-2}$ & $\dfrac{x+2}{(x-2)(x+2)}$\\[1.3em]
         $\dfrac{1}{x-2}+\dfrac{3}{x^{2}-4}$ & $\dfrac{x+2+3}{(x-2)(x+2)}$\\[1.3em]
         $\dfrac{x+2+3}{(x-2)(x+2)}$ & $\dfrac{x+5}{(x-2)(x+2)}$}}
\end{ealglossed}

%================================================================
\qtopic{\ealgloss{Multiplying}{ضرب دینا} and \ealgloss{dividing}{تقسیم کرنا} \ealgloss{algebraic}{الجبری} \ealgloss{fractions}{کسر}}
%================================================================

\begin{ealglossed}
\tf[\Large][\large]{To work out $\dfrac{a}{b}\div\dfrac{c}{d}$ you turn the
   first \ealgloss{fraction}{کسر} upside down}
  {\textbf{False.} It is the second \ealgloss{fraction}{کسر} that is turned upside down,
   and the two \ealgloss{fractions}{کسر} are then \ealgloss{multiplied}{ضرب دینا}.\\[0.3em]
   $\dfrac{a}{b}\div\dfrac{c}{d}=\dfrac{a}{b}\times\dfrac{d}{c}$}
\end{ealglossed}

\begin{ealglossed}
\tf[\Large][\large]{$\dfrac{x}{2}\times\dfrac{3}{5}$ can be worked out
   without finding a \ealgloss{common denominator}{مشترک مخرج}}
  {\textbf{True.} A \ealgloss{common denominator}{مشترک مخرج} is needed for \ealgloss{adding}{جمع کرنا} and
   \ealgloss{subtracting}{تفریق کرنا}, not for \ealgloss{multiplying}{ضرب دینا}.\\[0.3em]
   The \ealgloss{numerators}{صورت} are \ealgloss{multiplied}{ضرب دینا} and the \ealgloss{denominators}{مخرج} are \ealgloss{multiplied}{ضرب دینا}, so
   the answer is $\dfrac{3x}{10}$.}
\end{ealglossed}

\tf[\LARGE][\large]{$\dfrac{a}{b}\times\dfrac{c}{d}=\dfrac{a+c}{b+d}$}
  {\textbf{False.} Test it with $\dfrac{1}{2}\times\dfrac{1}{2}$: the left
   side is $\dfrac{1}{4}$ and the right side is $\dfrac{2}{4}=\dfrac{1}{2}$.\\[0.3em]
   $\dfrac{a}{b}\times\dfrac{c}{d}=\dfrac{ac}{bd}$}

\begin{ealglossed}
\tf[\Large][\large]{$\dfrac{x}{4}\div 2=\dfrac{x}{2}$}
  {\textbf{False.} The whole number $2$ is the \ealgloss{fraction}{کسر} $\dfrac{2}{1}$,
   so it is turned upside down and the two are \ealgloss{multiplied}{ضرب دینا}.\\[0.3em]
   $\dfrac{x}{4}\div\dfrac{2}{1}=\dfrac{x}{4}\times\dfrac{1}{2}=\dfrac{x}{8}$\\[0.3em]
   When $x=8$: $2\div 2=1$, not $4$.}
\end{ealglossed}

\begin{ealglossed}
\tf[\Large][\large]{$\dfrac{x+1}{4}\times\dfrac{8}{x+1}$ \ealgloss{simplifies}{سادہ کرنا} to $2$}
  {\textbf{True.} \ealgloss{Cancelling}{حذف کرنا} across a multiplication \ealgloss{sign}{علامت} is allowed,
   because $x+1$ is a \ealgloss{factor}{عامل} of a whole \ealgloss{numerator}{صورت} and of a whole
   \ealgloss{denominator}{مخرج}.\\[0.3em]
   $\dfrac{8\strike{(x+1)}}{4\strike{(x+1)}}=\dfrac{8}{4}=2$}
\end{ealglossed}

\begin{ealglossed}
\qq[\LARGE][\large]{\ealgloss{Simplify}{سادہ کرنا}\\[0.4em]
   $\dfrac{x}{2}\times\dfrac{3}{x}$}
  {\ealgloss{Multiply}{ضرب دینا} the \ealgloss{numerators}{صورت} and \ealgloss{multiply}{ضرب دینا} the \ealgloss{denominators}{مخرج}, then \ealgloss{cancel}{حذف کرنا} the
   \ealgloss{common factor}{مشترک عامل} $x$.
   \par\vspace{0.4em}
   \work{$\dfrac{x}{2}\times\dfrac{3}{x}$ & $\dfrac{3x}{2x}$\\[0.95em]
         $\dfrac{3\strike{x}}{2\strike{x}}$ & $\dfrac{3}{2}$}}
\end{ealglossed}

\begin{ealglossed}
\qq[\LARGE][\large]{\ealgloss{Simplify}{سادہ کرنا}\\[0.4em]
   $\dfrac{3x}{4}\times\dfrac{8}{9x}$}
  {\work{$\dfrac{3x}{4}\times\dfrac{8}{9x}$ & $\dfrac{24x}{36x}$\\[0.95em]
         $\dfrac{24\strike{x}}{36\strike{x}}$ & $\dfrac{24}{36}$\\[0.95em]
         $\dfrac{24}{36}$ & $\dfrac{2}{3}$}}
\end{ealglossed}

\begin{ealglossed}
\qq[\LARGE][\large]{\ealgloss{Simplify}{سادہ کرنا}\\[0.4em]
   $\dfrac{x}{5}\div\dfrac{x}{2}$}
  {\flow{\tfrac{x}{5}\div\tfrac{x}{2}}{turn the second \ealgloss{fraction}{کسر} upside
   down and \ealgloss{multiply}{ضرب دینا}}{\tfrac{x}{5}\times\tfrac{2}{x}}{\ealgloss{cancel}{حذف کرنا} the \ealgloss{common factor}{مشترک عامل} $x$}{\tfrac{2}{5}}}
\end{ealglossed}

\begin{ealglossed}
\qq[\LARGE][\large]{\ealgloss{Simplify}{سادہ کرنا}\\[0.4em]
   $\dfrac{x}{3}\div\dfrac{2}{x}$}
  {\flow{\tfrac{x}{3}\div\tfrac{2}{x}}{turn the second \ealgloss{fraction}{کسر} upside
   down and \ealgloss{multiply}{ضرب دینا}}{\tfrac{x}{3}\times\tfrac{x}{2}}{\ealgloss{multiply}{ضرب دینا} the
   \ealgloss{numerators}{صورت} and the \ealgloss{denominators}{مخرج}}{\tfrac{x^{2}}{6}}}
\end{ealglossed}

\begin{ealglossed}
\qq[\LARGE][\large]{\ealgloss{Simplify}{سادہ کرنا}\\[0.4em]
   $\dfrac{x+1}{3}\times\dfrac{6}{x+2}$}
  {The two \ealgloss{brackets}{بریکٹ} are different, so only the numbers \ealgloss{cancel}{حذف کرنا}.
   \par\vspace{0.4em}
   \work{$\dfrac{x+1}{3}\times\dfrac{6}{x+2}$ & $\dfrac{6(x+1)}{3(x+2)}$\\[0.95em]
         $\dfrac{6(x+1)}{3(x+2)}$ & $\dfrac{2(x+1)}{x+2}$}}
\end{ealglossed}

\begin{ealglossed}
\qq[\LARGE][\small]{\ealgloss{Simplify}{سادہ کرنا}\\[0.4em]
   $\dfrac{x+2}{5}\div\dfrac{x+2}{10}$}
  {Turn the second \ealgloss{fraction}{کسر} upside down and \ealgloss{multiply}{ضرب دینا}.
   \par\vspace{0.4em}
   \work{$\dfrac{x+2}{5}\div\dfrac{x+2}{10}$ & $\dfrac{x+2}{5}\times\dfrac{10}{x+2}$\\[0.95em]
         $\dfrac{x+2}{5}\times\dfrac{10}{x+2}$ & $\dfrac{10(x+2)}{5(x+2)}$\\[0.95em]
         $\dfrac{10\strike{(x+2)}}{5\strike{(x+2)}}$ & $2$}}
\end{ealglossed}

\begin{ealglossed}
\qq[\LARGE][\large]{\ealgloss{Simplify}{سادہ کرنا}\\[0.4em]
   $\dfrac{5x^{3}}{4}\div\dfrac{x}{8}$}
  {\flow{\tfrac{5x^{3}}{4}\div\tfrac{x}{8}}{turn the second \ealgloss{fraction}{کسر}
   upside down and \ealgloss{multiply}{ضرب دینا}}{\tfrac{5x^{3}}{4}\times\tfrac{8}{x}}{\ealgloss{divide}{تقسیم کرنا}
   the numbers and \ealgloss{subtract}{تفریق کرنا} the \ealgloss{indices}{اس}}{10x^{2}}}
\end{ealglossed}

\begin{ealglossed}
\qq[\Large][\small]{\ealgloss{Simplify}{سادہ کرنا}\\[0.4em]
   $\dfrac{x^{2}-9}{x+1}\times\dfrac{x+1}{x-3}$}
  {$x^{2}-9$ is the \ealgloss{difference of two squares}{دو مربعوں کا فرق}.
   \par\vspace{0.4em}
   \work{$x^{2}-9$ & $(x-3)(x+3)$\\[0.95em]
         $\dfrac{x^{2}-9}{x+1}\times\dfrac{x+1}{x-3}$ & $\dfrac{(x-3)(x+3)(x+1)}{(x+1)(x-3)}$\\[0.95em]
         $\dfrac{\strike{(x-3)}(x+3)\strike{(x+1)}}{\strike{(x+1)}\strike{(x-3)}}$ & $x+3$}}
\end{ealglossed}

\begin{ealglossed}
\qq[\Large][\small]{\ealgloss{Simplify}{سادہ کرنا}\\[0.4em]
   $\dfrac{x^{2}+7x+10}{x-1}\div\dfrac{x+2}{x-1}$}
  {\ealgloss{Factorise}{تجزیه کرنا} the \ealgloss{quadratic}{دو درجی}, then turn the second \ealgloss{fraction}{کسر} upside down and
   \ealgloss{multiply}{ضرب دینا}.
   \par\vspace{0.4em}
   \work{$x^{2}+7x+10$ & $(x+2)(x+5)$\\[0.95em]
         $\dfrac{(x+2)(x+5)}{x-1}\div\dfrac{x+2}{x-1}$ & $\dfrac{(x+2)(x+5)}{x-1}\times\dfrac{x-1}{x+2}$\\[0.95em]
         $\dfrac{\strike{(x+2)}(x+5)\strike{(x-1)}}{\strike{(x-1)}\strike{(x+2)}}$ & $x+5$}}
\end{ealglossed}

\begin{ealglossed}
\qq[\Large][\footnotesize]{\ealgloss{Simplify}{سادہ کرنا}\\[0.4em]
   $\dfrac{x^{2}+3x+2}{x^{2}-1}\div\dfrac{x^{2}+4x+4}{x^{2}-x}$}
  {\ealgloss{Factorise}{تجزیه کرنا} all four \ealgloss{quadratics}{دو درجی}, then turn the second \ealgloss{fraction}{کسر} upside
   down and \ealgloss{multiply}{ضرب دینا}.
   \par\vspace{0.4em}
   \work{$x^{2}+3x+2$ & $(x+1)(x+2)$\\[0.95em]
         $x^{2}-1$ & $(x+1)(x-1)$\\[0.95em]
         $x^{2}+4x+4$ & $(x+2)(x+2)$\\[0.95em]
         $x^{2}-x$ & $x(x-1)$\\[0.95em]
         $\dfrac{(x+1)(x+2)}{(x+1)(x-1)}\times\dfrac{x(x-1)}{(x+2)(x+2)}$ & $\dfrac{(x+1)(x+2)\,x\,(x-1)}{(x+1)(x-1)(x+2)(x+2)}$\\[1.5em]
         $\dfrac{\strike{(x+1)}\strike{(x+2)}\,x\,\strike{(x-1)}}{\strike{(x+1)}\strike{(x-1)}\strike{(x+2)}(x+2)}$ & $\dfrac{x}{x+2}$}}
\end{ealglossed}

\end{document}